% from ega4-13-fr.tex
\oldpage[IV-3]{196}
convenable de $A/\mathfrak p$, supposer que chaque $t_i$ est l'image canonique
dans $B\otimes_Ak(\mathfrak p)$ d'un élément $s_i\in B$. Soit alors
$u:A[T_1,\ldots,T_e]\to B$ l'homomorphisme tel que $u(T_i)=s_i$ pour tout
$i$, et soit $g:X\to Y[T_1,\ldots,T_e]$ le morphisme correspondant. Il est
clair qu'en raison du choix des $s_i$, $g_y$ est un morphisme fini. Pour tout
morphisme $g:X\to Y[T_1,\ldots,T_e]$, tel que $g_y$ soit fini, il résulte de
(5.4.2) et (4.1.2.1) que $g_y$ est nécessairement surjectif. D'autre part, en
vertu de (13.1.4), l'ensemble $U$ des $x\in X$ qui sont isolés dans leur fibre
$g^{-1}(g(x))$ est ouvert dans $X$ et contient $f^{-1}(y)$, et par définition
la restriction $g|U$ est un morphisme quasi-fini.

Ce lemme étant établi, pour prouver que $a$) implique $b$), il reste à voir
que si $x_\lambda$ est un point maximal de $U$, $g(x_\lambda)=z_\lambda$ est
un point maximal de $Z=Y[T_1,\ldots,T_e]$. En vertu de l'hypothèse $a$), on
peut (en restreignant au besoin $U$), supposer que $f(x_\lambda)$ est l'un des
points génériques $y_\alpha$ des composantes irréductibles de $Y$ contenant
$y$~; si $p:Z\to Y$ est le morphisme structural, on a donc
$p(z_\lambda)=y_\alpha$, et l'on déduit par suite de $g$ un
$k(y_\alpha)$-morphisme quasi-fini
\[
  h:U\cap f^{-1}(y_\alpha)\longrightarrow p^{-1}(y_\alpha).
\]
Or $p^{-1}(y_\alpha)=\operatorname{Spec}(k(y_\alpha)[T_1,\ldots,T_e])$ est
intègre et de dimension $e$~; si $z_\lambda$ n'était pas point maximal de
$Z$, il ne serait pas point maximal de $p^{-1}(y_\alpha)$
($\mathbf 0_{\mathrm I}$, 2.1.8) et son adhérence $Z'_\lambda$ dans
$p^{-1}(y_\alpha)$ serait donc de dimension $<e$ (4.1.2.1). Mais comme $h$
est quasi-fini, il résulte de (4.1.2) et de l'hypothèse $a$) que l'on a
$\dim(Z'_\lambda)\geq e$ (la restriction de $h$ à
$\overline{\{x_\lambda\}}$ se factorisant en
\[
  \overline{\{x_\lambda\}}\longrightarrow Z'_\lambda
  \longrightarrow p^{-1}(y_\alpha)
\]
en vertu de (\textbf{I}, 5.2.2))~; on aboutit ainsi à une contradiction, ce
qui montre que $a$) entraîne $b$).

Prouvons finalement que $b$) implique $a$). Notons que le morphisme structural
$p:Z=Y[T_1,\ldots,T_e]\to Y$ est fidèlement plat~; donc (2.3.4) les points
maximaux de $Z$ ont pour images par $p$ les points maximaux de $Y$~; ceci
prouve déjà que les $f(x_\lambda)$ sont points génériques des $Y_\alpha$. En
outre, si $f(x_\lambda)=y_\alpha$, le $k(y_\alpha)$-morphisme
$h:U\cap f^{-1}(y_\alpha)\to p^{-1}(y_\alpha)$ déduit de $g$ est dominant et
quasi-fini par hypothèse~; on a donc
$\dim(U_\lambda\cap f^{-1}(y_\alpha))=e$ en vertu de (4.1.2)
($U_\lambda$ étant la composante irréductible de $U$ de point générique
$x_\lambda$). De même, pour tout $x'\in U_\lambda$, le morphisme
$U_\lambda\cap f^{-1}(f(x'))\to p^{-1}(f(x'))$ déduit de $g$ est un
$k(f(x'))$-morphisme quasi-fini, donc (4.1.2) on a
$\dim(f^{-1}(f(x'))\cap U_\lambda)\leq e$~; mais par ailleurs on sait (13.1.6)
que toutes les composantes irréductibles de
$U_\lambda\cap f^{-1}(f(x'))$ sont de dimension $\geq e$~; on voit ainsi que
ces composantes sont exactement de dimension $e$, donc $b$) entraîne $a$).
C.Q.F.D.
\end{proof}

\begin{definition}[13.3.2]
\label{IV.13.3.2-fr}
Soient $Y$ un préschéma, $f:X\to Y$ un morphisme localement de type fini, $x$
un point de $X$. On dit que $f$ est équidimensionnel au point $x$ (ou que $X$
est équidimensionnel sur $Y$ au point $x$) si les conditions équivalentes de
(13.3.1) sont vérifiées. On dit que $f$ est équidimensionnel (ou que $X$ est
équidimensionnel sur $Y$) si $f$ est équidimensionnel en tout point $x\in X$.
\end{definition}

Il est clair que, pour que $f$ soit équidimensionnel en un point $x$, il faut
et il suffit que $f_{\mathrm{red}}$ le soit. En outre, les conditions de
(13.3.1) montrent que l'ensemble des points où $f$ est équidimensionnel est
\emph{ouvert} dans $X$.

On notera que lorsque $X$ et $Y$ sont \emph{irréductibles}, dire que $f$ est
équidimen-
